% ==============================================================================
%  4.4  El análisis de los residuos: gráficos y contrastes
% ==============================================================================
\section{El análisis de los residuos}
\label{sec:t4-analisis}

Las \textbf{representaciones gráficas} de los residuos son la principal herramienta de
validación. Las más utilizadas son:
\begin{itemize}
  \item residuos frente a la variable $X$;
  \item valor absoluto o cuadrado de los residuos frente a la variable $X$;
  \item residuos frente a los valores predichos;
  \item residuos frente al orden de las observaciones;
  \item residuos frente a otras posibles variables $X$;
  \item diagrama de cajas de los residuos;
  \item gráfico de normalidad de los residuos.
\end{itemize}
A continuación se indica qué gráficos permiten estudiar cada hipótesis del modelo.

\subsection{Linealidad}

La hipótesis de linealidad se estudia en el \textbf{gráfico de residuos}: residuos frente a $X$
o, equivalentemente, frente a los valores predichos. En el modelo de regresión lineal simple
ambos gráficos son totalmente equivalentes, ya que $\est{Y}_i=\est{\beta}_0+\est{\beta}_1X_i$ es una
transformación lineal de $X_i$ (con el eje horizontal invertido si $\est{\beta}_1<0$), por lo
que lo habitual es utilizar el de los valores predichos.

Cuando el modelo lineal es adecuado, los residuos se sitúan centrados en 0 de forma aleatoria,
sin que se observen patrones sistemáticos (\emph{gráfico nulo}). Si la función de regresión no es
lineal, los residuos presentan un patrón curvo.

\begin{figure}[H]
  \centering
  \begin{subfigure}[t]{0.32\linewidth}
    \centering
    \begin{tikzpicture}
      \begin{axis}[residuos, xlabel={$\est{y}_i$ o $x_i$}]
        \addplot[cero] {0};
        \addplot[puntos] table[x=x, y=r] {datos/tema4/nulo.dat};
      \end{axis}
    \end{tikzpicture}
    \caption{Gráfico nulo: modelo adecuado.}
  \end{subfigure}\hfill
  \begin{subfigure}[t]{0.32\linewidth}
    \centering
    \begin{tikzpicture}
      \begin{axis}[residuos]
        \addplot[cero] {0};
        \addplot[puntos] table[x=x, y=r] {datos/tema4/curva1.dat};
      \end{axis}
    \end{tikzpicture}
    \caption{Patrón curvo: relación no lineal.}
  \end{subfigure}\hfill
  \begin{subfigure}[t]{0.32\linewidth}
    \centering
    \begin{tikzpicture}
      \begin{axis}[residuos]
        \addplot[cero] {0};
        \addplot[puntos] table[x=x, y=r] {datos/tema4/curva2.dat};
      \end{axis}
    \end{tikzpicture}
    \caption{Patrón curvo: relación no lineal.}
  \end{subfigure}
  \caption{Gráficos de residuos y linealidad (datos simulados).}
\end{figure}

\subsection{Homogeneidad de varianzas}

La hipótesis de homogeneidad de varianzas también se estudia en el gráfico de residuos. Si la
varianza es homogénea, la dispersión de los residuos es aproximadamente constante a lo largo de
todo el eje horizontal; en caso contrario, la nube de puntos se abre o se cierra.

\begin{figure}[H]
  \centering
  \begin{subfigure}[t]{0.45\linewidth}
    \centering
    \begin{tikzpicture}
      \begin{axis}[residuos, xlabel={$\est{y}_i$ o $x_i$}, ymin=-3, ymax=3]
        \addplot[cero] {0};
        \addplot[puntos] table[x=x, y=r] {datos/tema4/nulo.dat};
      \end{axis}
    \end{tikzpicture}
    \caption{Varianza constante.}
  \end{subfigure}\hfill
  \begin{subfigure}[t]{0.45\linewidth}
    \centering
    \begin{tikzpicture}
      \begin{axis}[residuos, xlabel={$\est{y}_i$ o $x_i$}]
        \addplot[cero] {0};
        \addplot[puntos] table[x=x, y=r] {datos/tema4/embudo_crec.dat};
      \end{axis}
    \end{tikzpicture}
    \caption{La varianza aumenta con $\est{y}_i$.}
  \end{subfigure}

  \medskip
  \begin{subfigure}[t]{0.45\linewidth}
    \centering
    \begin{tikzpicture}
      \begin{axis}[residuos, xlabel={$\est{y}_i$ o $x_i$}]
        \addplot[cero] {0};
        \addplot[puntos] table[x=x, y=r] {datos/tema4/embudo_decr.dat};
      \end{axis}
    \end{tikzpicture}
    \caption{La varianza disminuye con $\est{y}_i$.}
  \end{subfigure}\hfill
  \begin{subfigure}[t]{0.45\linewidth}
    \centering
    \begin{tikzpicture}
      \begin{axis}[residuos, xlabel={$\est{y}_i$ o $x_i$}]
        \addplot[cero] {0};
        \addplot[puntos] table[x=x, y=r] {datos/tema4/diamante.dat};
      \end{axis}
    \end{tikzpicture}
    \caption{Mayor varianza en los valores centrales.}
  \end{subfigure}
  \caption{Gráficos de residuos y homogeneidad de varianzas (datos simulados).}
\end{figure}

También son útiles los gráficos del valor absoluto o del cuadrado de los residuos frente a
$\est{Y}_i$ (o $X_i$), especialmente cuando no hay muchas observaciones: si la varianza no es
homogénea, la magnitud de los residuos, independientemente de su signo, cambia con el valor de
$\est{Y}_i$ (o $X_i$).

\begin{nota}[Residuos regulares frente a estudentizados]
  Como $\Var(e_i)=\sigma^2(1-h_{ii})$ (\cref{prop:t4-prop-residuos}) y $h_{ii}$ es mínimo cuando
  $X_i=\media{X}$, la varianza de los residuos regulares es \emph{mayor} cuanto más cerca está
  $X_i$ de su media muestral.%
   Este patrón podría verse en el gráfico de residuos aunque los errores
  tengan varianza constante. Por ese motivo es mucho más frecuente utilizar los residuos
  estudentizados, que tienen varianza 1 para todas las observaciones.
\end{nota}

\subsection{Normalidad}

Para estudiar la normalidad de los errores se utilizan:
\begin{itemize}
  \item el \textbf{gráfico cuantil-cuantil} (\emph{QQ-plot}), en el que los residuos ordenados se
    comparan con sus valores esperados bajo normalidad, es decir, con los cuantiles de la normal.
    Si los residuos son normales, los puntos se sitúan aproximadamente sobre una recta;
  \item el \textbf{histograma} o el \textbf{diagrama de cajas} de los residuos.
\end{itemize}

\begin{figure}[H]
  \centering
  \begin{subfigure}[t]{0.32\linewidth}
    \centering
    \begin{tikzpicture}
      \begin{axis}[mini, xlabel={$z$}, ylabel={$r_i$}, xmin=-3, xmax=3, ymin=-4, ymax=6]
        \addplot[cero, domain=-3:3] {x};
        \addplot[puntos] table[x=q, y=der] {datos/tema4/qq.dat};
      \end{axis}
    \end{tikzpicture}
    \caption{Distribución asimétrica a la derecha.}
  \end{subfigure}\hfill
  \begin{subfigure}[t]{0.32\linewidth}
    \centering
    \begin{tikzpicture}
      \begin{axis}[mini, xlabel={$z$}, ylabel={$r_i$}, xmin=-3, xmax=3, ymin=-6, ymax=4]
        \addplot[cero, domain=-3:3] {x};
        \addplot[puntos] table[x=q, y=izq] {datos/tema4/qq.dat};
      \end{axis}
    \end{tikzpicture}
    \caption{Distribución asimétrica a la izquierda.}
  \end{subfigure}\hfill
  \begin{subfigure}[t]{0.32\linewidth}
    \centering
    \begin{tikzpicture}
      \begin{axis}[mini, xlabel={$z$}, ylabel={$r_i$}, xmin=-3, xmax=3, ymin=-5, ymax=5]
        \addplot[cero, domain=-3:3] {x};
        \addplot[puntos] table[x=q, y=colas] {datos/tema4/qq.dat};
      \end{axis}
    \end{tikzpicture}
    \caption{Distribución simétrica con colas pesadas.}
  \end{subfigure}
  \caption{QQ-plots de residuos no normales: residuos ordenados ($r_i$) frente a los cuantiles
    esperados de la normal ($z$). La línea discontinua es la recta que seguirían los puntos bajo
    normalidad (datos simulados).}
\end{figure}

Con un tamaño muestral adecuado también es útil comparar las frecuencias observadas con las
esperadas bajo normalidad: aproximadamente el 68\,\% de los residuos deben quedar en la banda
$\pm\sqrt{MSE}$, el 95\,\% en $\pm2\sqrt{MSE}$ y el 99.7\,\% en $\pm3\sqrt{MSE}$.

Para muestras grandes, el supuesto de normalidad no es crucial: \textbf{las inferencias sobre los
parámetros de la regresión $\beta_0$ y $\beta_1$ son robustas ante la falta de normalidad de los
errores}, ya que con tamaños muestrales suficientes la aproximación a la normal es razonable por
el teorema central del límite (\cref{sec:t3-coeficientes}).

\subsection{Independencia}

Si las observaciones proceden de una muestra aleatoria de sujetos de alguna población, en
principio serán independientes. Para detectar correlación entre los errores se representan los
residuos frente al tiempo o frente a algún otro tipo de secuencia (\textbf{gráficos
secuenciales}). La presencia de \textbf{rachas} en los residuos sugiere falta de independencia.

\begin{figure}[H]
  \centering
  \begin{subfigure}[t]{0.48\linewidth}
    \centering
    \begin{tikzpicture}
      \begin{axis}[width=6.2cm, height=4.4cm, xlabel={Orden temporal}, ylabel={$e_i$},
          xmin=0, xmax=14, ymin=-2.6, ymax=2.4, ytick={0}, xtick={1,3,...,13},
          label style={font=\small}, tick label style={font=\footnotesize},
          ylabel style={rotate=-90}]
        \addplot[black, no marks, domain=0:14] {0};
        \addplot[puntos, mark size=1.3pt] coordinates {(1,-2.2) (2,-1.4) (3,-1.6) (4,-1.2)
          (5,-1.0) (6,-0.7) (7,0.3) (8,-0.1) (9,0.9) (10,1.3) (11,1.6) (12,1.3) (13,1.8)};
      \end{axis}
    \end{tikzpicture}
    \caption{Efecto de tendencia.}
  \end{subfigure}\hfill
  \begin{subfigure}[t]{0.48\linewidth}
    \centering
    \begin{tikzpicture}
      \begin{axis}[width=6.2cm, height=4.4cm, xlabel={Orden temporal}, ylabel={$e_i$},
          xmin=0, xmax=15, ymin=-2.2, ymax=2.4, ytick={0}, xtick={1,3,...,13},
          label style={font=\small}, tick label style={font=\footnotesize},
          ylabel style={rotate=-90}]
        \addplot[black, no marks, domain=0:15] {0};
        \addplot[puntos, mark size=1.3pt] coordinates {(1,-1.5) (2,-1.3) (3,-0.6) (4,0.2)
          (5,0.15) (6,1.0) (7,1.8) (8,1.2) (9,0.2) (10,0.15) (11,-1.1) (12,-1.4) (13,-1.2)
          (14,-0.3)};
      \end{axis}
    \end{tikzpicture}
    \caption{Falta de independencia cíclica.}
  \end{subfigure}
  \caption{Gráficos secuenciales de residuos con rachas (datos reconstruidos a partir de las
    diapositivas).}
\end{figure}

En el gráfico de residuos frente a $X$ la nube de puntos puede parecer no aleatoria sin que
exista falta de independencia. Normalmente, en ese caso, el problema es que $X$ no es un buen
predictor de $Y$.

\subsection{Identificación de variables \texorpdfstring{$X$}{X} importantes}
\label{sec:t4-omitidas}

Representar los residuos de un modelo frente a variables $X$ que no se han incluido en él puede
indicar que dicha variable tiene sobre la respuesta un efecto importante que el modelo no
recoge.

\begin{ejemplo}[Masa muscular y edad][ej:t4-masa]
  Se espera que la masa muscular de una persona decrezca con la edad. Con $Y$ = masa muscular y
  $X$ = edad se ajusta el modelo de regresión lineal simple
  \[
    Y_i = \beta_0+\beta_1X_i+\varepsilon_i, \qquad \varepsilon_i\iid\Normal(0,\sigma^2), \quad i=1,\ldots,n,
  \]
  a una muestra de hombres y mujeres, y se representan los residuos estudentizados eliminados
  frente a la edad, distinguiendo el sexo de cada individuo.
  \begin{center}
    \begin{tikzpicture}
      \begin{axis}[width=9cm, height=5.5cm, xlabel={Edad (años)}, ylabel={$t_{(i)}$},
          xmin=40, xmax=80, ymin=-3.5, ymax=3.5, ytick={-3,-2,...,3},
          label style={font=\small}, tick label style={font=\footnotesize},
          ylabel style={rotate=-90},
          legend style={font=\scriptsize, at={(0.5,1.03)}, anchor=south, draw=none,
            legend columns=2}]
        \addplot[only marks, mark=*, mark size=1.1pt, NavyBlue]
          table[x=edad, y=t, restrict expr to domain={\thisrow{sexo}}{0:0}] {datos/tema4/masa.dat};
        \addlegendentry{Mujeres}
        \addplot[only marks, mark=triangle*, mark size=1.4pt, Maroon]
          table[x=edad, y=t, restrict expr to domain={\thisrow{sexo}}{1:1}] {datos/tema4/masa.dat};
        \addlegendentry{Hombres}
        \addplot[gray, no marks, domain=40:80, forget plot] {0};
        \addplot[gray, dashed, no marks, domain=40:80, forget plot] {3};
        \addplot[gray, dashed, no marks, domain=40:80, forget plot] {-3};
      \end{axis}
    \end{tikzpicture}\\
    {\footnotesize (Datos simulados con el patrón del gráfico de las diapositivas.)}
  \end{center}
  Los residuos de los hombres son mayoritariamente positivos y los de las mujeres
  mayoritariamente negativos: el modelo subestima la masa muscular de los hombres y la
  sobrestima en las mujeres. El sexo tiene, por tanto, un efecto importante sobre la masa
  muscular que el modelo no recoge, y convendría incluirlo como variable explicativa.
\end{ejemplo}

\subsection{Contrastes de hipótesis sobre los residuos}

Si a la vista de los gráficos no está claro qué concluir, pueden utilizarse ciertos
\textbf{contrastes de hipótesis}. Estos contrastes \textbf{no reemplazan a las herramientas
gráficas}, ya que dependen mucho del tamaño muestral, sino que se utilizan como complemento. Los
más habituales son:
\begin{itemize}
  \item \emph{Homogeneidad de varianzas}: test de Brown--Forsythe y test de Breusch--Pagan.
  \item \emph{Normalidad}: test de Shapiro--Wilk y test de Kolmogorov--Smirnov, entre otros.
  \item \emph{Independencia}: test de Durbin--Watson.
\end{itemize}

\subsubsection{Contraste de homogeneidad de varianzas: test de Brown--Forsythe}

El más habitual es el test de Brown--Forsythe, una modificación del test de Levene que es robusta
frente a la falta de normalidad.
\begin{itemize}
  \item La muestra se divide en dos grupos: valores bajos y valores altos de $X$.
  \item Si la varianza no es constante, los residuos de un grupo serán más variables que los del
    otro, y las desviaciones absolutas de los residuos alrededor del centro de su grupo serán
    distintas en ambos grupos.
  \item El test se basa en el contraste $t$ de Student para dos muestras aplicado a esas
    desviaciones.
\end{itemize}
Sean $e_{i1}$, $i=1,\ldots,n_1$, y $e_{i2}$, $i=1,\ldots,n_2$, los residuos de los grupos 1 y 2
($n_1+n_2=n$), y $\tilde e_1$, $\tilde e_2$ sus medianas. Se definen
\[
  d_{i1} = |e_{i1}-\tilde e_1|, \qquad d_{i2} = |e_{i2}-\tilde e_2| ,
\]
con medias $\media{d}_1$ y $\media{d}_2$. El estadístico es
\[
  T_{BF} = \frac{\media{d}_1-\media{d}_2}{S\sqrt{\dfrac{1}{n_1}+\dfrac{1}{n_2}}},
  \qquad
  S^2 = \frac{\sum_{i=1}^{n_1}\big(d_{i1}-\media{d}_1\big)^2+\sum_{i=1}^{n_2}\big(d_{i2}-\media{d}_2\big)^2}{n_1+n_2-2},
\]
que, bajo la hipótesis de varianza constante, sigue aproximadamente una $\tdist{n_1+n_2-2}$. Se
rechaza la homogeneidad de varianzas si $|T_{BF}|$ es grande comparado con los cuantiles de esa
distribución.

\subsubsection{Contrastes de normalidad}

En general, estos contrastes sirven para contrastar hipótesis sobre la distribución de una
variable, no solo la normalidad. Se aplican a los residuos; a continuación se describen para
una muestra genérica $x_1,\ldots,x_n$.

\paragraph{Test $\chi^2$ de Pearson.} Es un contraste de bondad de ajuste a una distribución $F$:
\[
  H_0\colon \text{la distribución poblacional es } F
  \quad\text{frente a}\quad
  H_1\colon \text{la distribución poblacional no es } F .
\]
Se basa en el estadístico
\[
  Q = \sum_{j=1}^k \frac{(O_j-E_j)^2}{E_j} \bajoHo \chisq{k-1} \ \text{(aproximadamente)},
\]
donde $k$ es el número de niveles de una variable discreta o de clases de una continua, y $O_j$ y
$E_j$ son las frecuencias observadas y esperadas bajo $F$ en el nivel o clase $j$. Los $k-1$
grados de libertad corresponden a una $F$ completamente especificada; si se estiman $m$
parámetros de $F$ a partir de la muestra (por ejemplo, la varianza de los errores), son $k-1-m$.

\paragraph{Test de Kolmogorov--Smirnov.} Es un contraste de bondad de ajuste a distribuciones
continuas. Se basa en la máxima distancia entre la función de distribución empírica de la
muestra, $F_n$, y la función de distribución teórica $F$:
\[
  D_n = \sup_x\big|F_n(x)-F(x)\big| .
\]
Valores grandes de $D_n$ llevan a rechazar $H_0$.

\begin{definicion}[Función de distribución empírica][def:t4-empirica]
  Sean $x_{(1)}\le\cdots\le x_{(n)}$ los datos ordenados de menor a mayor. La \textbf{función de
  distribución empírica} de la muestra es
  \[
    F_n(x) =
    \begin{cases}
      0 & \text{si } x < x_{(1)}, \\[2pt]
      \dfrac{\operatorname{card}\{j : x_j\le x\}}{n} & \text{si } x_{(i)}\le x<x_{(i+1)},\quad i=1,\ldots,n-1, \\[8pt]
      1 & \text{si } x\ge x_{(n)},
    \end{cases}
  \]
  donde $\operatorname{card}\{j : x_j\le x\}$ es el número de observaciones menores o iguales
  que $x$. Es una función escalonada y no decreciente, que en cada observación da un salto igual
  a la proporción de observaciones de la muestra que toman ese valor.
\end{definicion}

\paragraph{Test de Shapiro--Wilk.} Es útil con muestras pequeñas, ya que es más potente que el de
Kolmogorov--Smirnov. Se basa en un estadístico que cuantifica las diferencias entre las
observaciones y los valores esperados bajo la distribución normal. Toma valores entre 0 y 1, con
valores pequeños para muestras alejadas de la normalidad.

\paragraph{Test de asimetría.}
En las distribuciones simétricas, como la normal, el coeficiente de
asimetría poblacional es 0. Con los momentos centrales muestrales
$m_k=\frac1n\sumi (x_i-\media{x})^k$ y $S_X^2=m_2$, el coeficiente de asimetría se estima por
\[
  CA = \frac{m_3}{S_X^3} = \frac{\sumi (x_i-\media{x})^3}{nS_X^3}
  \ \bajoHo\ \Normal\!\left(0,\ \frac6n\right) \ \text{(aproximadamente)},
\]
o, equivalentemente, estandarizando, $CAS=\sqrt{\frac n6}\,CA\bajoHo\Normal(0,1)$.

\paragraph{Test de apuntamiento.} La distribución normal se caracteriza porque su coeficiente de
apuntamiento es 0. Este se estima por
\[
  CAp = \frac{m_4}{S_X^4}-3 = \frac{\sumi (x_i-\media{x})^4}{nS_X^4}-3
  \ \bajoHo\ \Normal\!\left(0,\ \frac{24}{n}\right) \ \text{(aproximadamente)},
\]
o, equivalentemente, estandarizando, $CApS=\sqrt{\frac{n}{24}}\,CAp\bajoHo\Normal(0,1)$.

En ambos casos se rechaza la normalidad si el estadístico estandarizado es grande en valor
absoluto comparado con los cuantiles de la $\Normal(0,1)$.
